% Source-supported supplement. Included by the main editor after the modeling tutorials.
% No new GUI or numerical execution is claimed by the instructions below.
\chapter{建模工具详解}\label{ch:modeling-tools}\label{sec:modeling-reference}

前两章通过具体模型介绍常见的建模过程。本章按操作目标说明各项工具，便于在编辑时查阅参数和检查方法。第一次使用某一组工具时，先用小规模模型完成一个操作，确认作用范围与结果，再把它用于自己的材料。选择、晶胞处理、软物质构建和过程记录各有相应入口，可以从正在处理的问题开始阅读。

建模页面随当前文档改变。分子中的变换和测量主要在\ui{分子建模}的\ui{原子编辑}、\ui{几何编辑}中；晶体中的局部分子使用\ui{晶体建模}的\ui{组分变换}、\ui{组分几何}。改变晶体整体时则使用\ui{晶格编辑}或\ui{超胞与晶格}。命令呈灰色时，先检查当前文档类型和选择数量。例如，\ui{移动原子}要求一个独立位点，\ui{设置键角}要求依次指定三个原子。

\section{精确编辑结构}

\subsection{保存并再次调用一组原子}\label{sec:selection-set-workflow}
\index{Xuanzeji@选择集}

选择集给一组模型位点起一个名字，适合反复处理表面底层、吸附分子或某一段链。它保存的是位点的身份，既不是屏幕上的一圈球，也不是另存的分子副本。

\begin{steps}
  \item 在三维视图中选中目标。单击选一个原子，Shift 或 Ctrl 配合单击可调整多选；\ui{B}进入框选，\ui{L}进入套索。若要选择整个分子，可双击一个原子；也可在原子右键菜单中选择\ui{Select Connected Component}。连通关系取决于当前显示模型使用的键，晶体网络可能因此选中很大一片结构。
  \item 若要选择一种元素，在该元素的原子上打开右键菜单，选择\ui{Select All \ldots\ Atoms}。先查看选择数，再开始替换或删除。
  \item 打开顶部建模页面的\ui{选择集}，单击\ui{保存当前选择}，输入能说明用途的名称，例如“底层固定原子”。
  \item 单击空白区域或按 Esc 清除选择，再从\ui{选择集}菜单中单击刚保存的名称，确认相同位点重新亮起。
  \item 保存项目。若还要保留这组选择的完整模型信息，可同时导出原生结构 JSON；CIF、XYZ 等交换格式并不承担保存全部选择信息的用途。
\end{steps}

菜单名称后的两个数表示当前还能找到的位点数与原先保存的位点数。例如，某组中的原子被删除后，前一个数会变小；召回时也会提示有位点已不存在。不能把剩余选择当成最初的完整一组。对同名选择重新执行保存，会更新该组内容；当前菜单提供保存与召回，没有独立的选择集删除按钮。

周期结构应额外核对\ui{sites}。显示重复范围改变后，同一位点可能出现多个映像；选择集仍按模型位点工作，不会因为屏幕上的球变多而自动增加成员。

\subsection{修正原子的种类与位置}\label{sec:atom-edit-reference}

\ui{添加原子}直接加入一个位点。晶体可以选择分数坐标或笛卡尔坐标，分子使用以\AA 为单位的笛卡尔坐标。与之不同，分子中的\ui{绘制原子}还可以同时指定连接对象与键级；要建立化学连接，应使用草绘或添加键，而不是只把两个原子放得很近。

\begin{steps}
  \item 在\ui{原子编辑}中打开\ui{添加原子}，输入元素，确认坐标类型，再分别填写三个坐标分量。先预览，确认没有放到已有原子上，再应用。
  \item 要替换已有原子，先选择它们，再打开\ui{替换元素}，填写目标元素。替换会保留这些位点的位置，不会重新搜索适合该元素的几何结构。
  \item 要准确移动一个位点，打开\ui{移动原子}，输入该原子的目标位置。晶体中先切换坐标类型再输入，避免把分数坐标当成\AA。
  \item 要让一组原子一起移动，使用\ui{平移原子}。例如，笛卡尔向量 \((0,0,1)\) 表示整体沿全局 $z$ 轴移动 $1\,\text{\AA}$；晶体中的分数向量 \((0,0,1)\) 表示移动一个 $\mathbf c$ 晶格矢量。
  \item 删除时先核对独立位点数，打开\ui{删除原子}并检查序号。应用后检查组成和剩余连接，最后另存结构或项目。
\end{steps}

这些局部晶体命令还可能出现\ui{保留对称性}或\ui{打破对称性}的选择。保持对称性会按等价关系处理相关位点；做局部掺杂或缺陷时，应选择与实际目的相符的方式。执行前看预览中的所有变化，不能只观察最初点中的那一个原子。

\subsection{整理重复位点和原子顺序}\label{sec:atom-clean-reference}

\ui{合并原子}处理相距非常近的晶体位点。它按以\AA 为单位的容差比较周期距离，不是把重复显示的边界映像变成更少的球。默认容差为 $0.01\,\text{\AA}$；初次处理导入模型时，应先用较小的容差检查。

\begin{steps}
  \item 保存原始模型副本，打开\ui{原子编辑}中的\ui{合并原子}。
  \item 选择合并方式。\ui{删除重合原子}保留一份组分信息；\ui{占位率相加}会合并同簇位点的物种占位；\ui{属性取平均}用于可兼容的数值位点属性。各方式都会把同簇位置移到按邻近周期映像求得的中心，不能理解为坐标一定保留不变。
  \item 预览后检查位点数、占位和局部位置。若两个真实相邻原子被并成一个，取消并减小容差。
  \item 应用后保存为新名称，保留原输入供比较。
\end{steps}

\ui{原子排序}按元素重新排列位点，\ui{倒序排列}改变排序方向。排序用于整理文件和序号，不改变原子的空间位置。已经记下的“第几个原子”应重新核对；需要长期追踪的一组位点更适合用选择集保存。

\ui{扰动原子}和\ui{扰动结构}提供随机几何扰动，参数是最小、最大位移，单位均为\AA，以及随机种子。例如，可用最小位移 0、最大位移 $0.05\,\text{\AA}$、种子 7 做一次很小的扰动练习。当前这两个命令作用于整个模型，不能把事先选中的少量原子理解为唯一扰动范围。预览时核对未选中的原子也可能移动；需要原始对照时，先保存副本。

\subsection{改变分子的方向}\label{sec:orientation-reference}

居中、旋转和镜像用于改变几何摆放。它们本身不执行能量优化。晶体中的吸附物也可使用\ui{组分变换}进行这些操作，但应先选择完整组分。

\subsubsection{把一组原子移到目标中心}

居中时，同一组中的所有原子使用同一个位移。设选择中有 $N$ 个原子，分子用笛卡尔位置 $\mathbf r_i$，晶体用分数坐标 $\mathbf f_i$，两种情况分别为
\begin{equation}
 \begin{aligned}
 \Delta\mathbf r&=\mathbf r_{\mathrm{target}}-
        \frac1N\sum_{i=1}^{N}\mathbf r_i,
 &\mathbf r_i'&=\mathbf r_i+\Delta\mathbf r,\\
 \Delta\mathbf f&=\mathbf f_{\mathrm{target}}-
        \frac1N\sum_{i=1}^{N}\mathbf f_i,
 &\mathbf f_i'&=\mathbf f_i+\Delta\mathbf f.
 \end{aligned}
 \label{eq:center-selected-atoms}
\end{equation}
第一行用于分子，位置与位移的单位都是\AA；第二行用于晶体，各量均为无量纲分数坐标。分子常用目标 $(0,0,0)$，晶体常用 $(0.5,0.5,0.5)$。每个原子在平均值中占相同权重，因此居中不使用原子质量，也不改变所选原子之间的距离。

例如，分子中两个原子的 $x$ 坐标为1和3\,\AA，要把它们的平均位置移到原点，共同位移就是 $-2\,\text{\AA}$；处理后为 $-1$ 和1\,\AA，间距仍为2\,\AA。晶体中若两个位点的第一分数分量为0.1和0.3，目标为0.5，则都加0.3，得到0.4和0.6。这里只取一个分量说明计算，其余两个分量同样处理。

周期边界附近要先看清使用的是哪一个映像。0.99与0.01的普通平均为0.5，但两个位点可能实际聚集在边界两侧；若要以这段连续结构的中心为准，应先整理为相邻的连续映像再居中。应用后坐标可能被归回晶胞，因此还应从周期连接看整体，而不只比较某个坐标的正负。

\subsubsection{绕指定锚点旋转}

旋转先从每个位置中减去锚点，再转动这个相对向量，最后加回锚点。这样，锚点本身不动，离锚点的距离也不变。沿用行向量表示，有
\begin{equation}
 \mathbf r'=\mathbf p+(\mathbf r-\mathbf p)
              R(\hat{\mathbf n},\theta)^{\mathsf T}.
 \label{eq:rotate-row-position}
\end{equation}
$\mathbf p$ 是以\AA 为单位的锚点，$\hat{\mathbf n}$ 是旋转轴的单位方向，$\theta$ 是转角；上标 $\mathsf T$ 表示转置。对话框中填写度数，正角遵循右手定则：右手拇指沿输入轴的正向，四指弯曲的方向为正转向。轴向量只决定方向，例如\literalinput{0 0 1}和\literalinput{0 0 2}表示同一根正向转轴。

绕全局 $z$ 轴旋转时，常用的矩阵为
\[
 R_z(\theta)=\begin{bmatrix}
 \cos\theta&-\sin\theta&0\\
 \sin\theta&\cos\theta&0\\
 0&0&1
 \end{bmatrix}.
\]
取锚点 $\mathbf p=(1,1,0)\,\text{\AA}$，位置 $\mathbf r=(2,1,0)\,\text{\AA}$，绕 $z$ 轴正转 $90^\circ$，得到 $\mathbf r'=(1,2,0)\,\text{\AA}$。这个小例子可用来辨认旋转方向；旋转后原子到锚点仍相距1\,\AA。晶体中原子的等价位置还会归回晶胞，预览时应连同边界映像一起观察。

\subsubsection{关于一个平面作镜像}

镜像需要一个平面上的点 $\mathbf p$ 和非零笛卡尔法向 $\mathbf n$。先把法向归一化，再减去两倍的有向投影：
\begin{equation}
 \begin{aligned}
 \hat{\mathbf n}&=\frac{\mathbf n}{\|\mathbf n\|},\\
 \mathbf r'&=\mathbf r-
 2\bigl[(\mathbf r-\mathbf p)\cdot\hat{\mathbf n}\bigr]
       \hat{\mathbf n}.
 \end{aligned}
 \label{eq:mirror-cartesian-plane}
\end{equation}
点积 $(\mathbf r-\mathbf p)\cdot\hat{\mathbf n}$ 是原子到平面的有向距离，单位为\AA。镜像后它的符号反转，大小相同；位于平面上的原子保持不动。法向填\literalinput{1 0 0}，平面点取 $(1,0,0)\,\text{\AA}$，就是 $x=1\,\text{\AA}$ 的平面。位置 $(1.5,2,0)\,\text{\AA}$ 镜像后变为 $(0.5,2,0)\,\text{\AA}$，两侧到平面均为0.5\,\AA。

\begin{figure}[!htb]
\centering
\begin{tikzpicture}[x=1mm,y=1mm,font=\sffamily\small]
 \node[text=brand] at (24,47) {绕锚点旋转};
 \draw[teal,dashed] (10,15)--(33,15);
 \draw[teal,dashed] (10,15)--(10,38);
 \draw[->,>=stealth,brand,line width=.85pt] (33,15) arc[start angle=0,end angle=90,radius=23];
 \fill[ink] (10,15) circle (1.1);
 \fill[teal] (33,15) circle (1.4);
 \fill[brand] (10,38) circle (1.4);
 \node[anchor=north] at (10,13) {$\mathbf p$};
 \node[anchor=west,text=teal] at (35,15) {$\mathbf r$};
 \node[anchor=east,text=brand] at (8,38) {$\mathbf r'$};
 \node[text=brand] at (30,35) {$+90^\circ$};
 \node[text=muted,font=\sffamily\footnotesize] at (24,3) {到锚点的距离保持不变};
 \node[text=brand] at (86,47) {关于平面镜像};
 \fill[paleblue] (83,8) rectangle (87,40);
 \draw[teal,line width=.8pt] (85,8)--(85,40);
 \draw[dashed,muted] (68,29)--(102,29);
 \draw[->,>=stealth,brand,line width=.7pt] (85,12)--(105,12) node[above] {$\hat{\mathbf n}$};
 \fill[ink] (85,18) circle (1.1);
 \fill[teal] (102,29) circle (1.4);
 \fill[brand] (68,29) circle (1.4);
 \node[anchor=south,text=teal] at (102,31) {$\mathbf r$};
 \node[anchor=south,text=brand] at (68,31) {$\mathbf r'$};
 \node[anchor=east] at (83,18) {$\mathbf p$};
 \node[text=muted,font=\sffamily\footnotesize] at (86,3) {到镜面的距离两侧相等};
\end{tikzpicture}
\caption{锚点旋转与平面镜像（截面示意）。左图中的点绕 $\mathbf p$ 旋转；右图中的点跨过镜面，沿法向到镜面的距离大小不变。}
\label{fig:round10-anchor-mirror}
\end{figure}


上述三个例子只说明参数的几何关系，不代表一次能量优化。实际处理分子时，先确认所选原子构成需要整体移动的部分，再按下面的步骤选择一种操作；不要把相机转动当作结构已经旋转。


\begin{steps}
  \item 选择要处理的一组原子，打开\ui{原子居中}。分子的目标中心以\AA 表示，常用 \((0,0,0)\)；晶体的目标中心为分数坐标，常用 \((0.5,0.5,0.5)\)。这里的中心是所选原子位置的算术平均，不是质量中心。
  \item 要作已知角度的转动，打开\ui{旋转原子}。角度用度，旋转轴为全局笛卡尔方向，锚点用\AA。例如，轴 \((0,0,1)\)、锚点 \((0,0,0)\)、角度 $90^\circ$ 表示绕过原点的 $z$ 轴转动。
  \item 要镜像，打开\ui{镜像原子}。填写平面法向和一个平面内点。例如，法向 \((1,0,0)\)、平面点 \((0,0,0)\) 定义 $yz$ 平面。法向只确定方向，不是一个需要输入\AA 的长度。
  \item 预览时比较原子间距离、连接与附近宿主。选完整分子作刚性变换时，内部分子几何应保持不变；只选一部分则可能拉断或扭曲原有几何。
  \item 应用并另存。若只是调整观看方向，退出建模工具后旋转相机即可，无须改变结构坐标。
\end{steps}

\ui{几何对齐}适合把长分子的主轴放到指定方向，或让一组近似共面的原子的平面法向对齐。选择\ui{principal\_axis}或\ui{plane\_normal}，再输入非零目标向量。前者看整体主方向，后者看拟合平面；对高度对称的分子，主轴可能不是唯一的，完成后仍应从不同方向查看。

\subsection{修改显式化学键}\label{sec:bond-edit-reference}

显式化学键记录哪两个原子相连及其键级；三维显示中的线条还可能来自按距离推断的邻接关系。修改拓扑时，应以分子中的连接记录为准。键级 1、2、3 分别表示单、双、三键，1.5 用于这里的芳香连接。

\begin{steps}
  \item 导入或打开要修改的分子，保存一份副本。依次选择构成目标键的两个原子，查看序号，不能多选。
  \item 两原子尚未连接时，在\ui{分子建模}菜单中打开\ui{添加化学键}，指定键级后应用。
  \item 已有连接需要改变键级时，使用\ui{设置键级}；需要断开连接时，使用\ui{删除化学键}。删除键保留两个原子，删除原子则会改变原子数，两者应按目的区分。
  \item 在\ui{原子化学属性}中核对受影响原子的元素、形式电荷、杂化与芳香标记。不替换元素时保留\ui{unchanged}。形式电荷是选定原子的属性，不是把整个分子的净电荷平均分配到所有原子。
  \item 若原有氢已不适合修改后的连接，选择相应重原子后先用\ui{删除氢}，再用\ui{添加氢}按新连接补齐。核对组成、重原子连接和氢数，并测量可疑的短距离；必要时再整理局部几何。
  \item 保存项目与原生结构 JSON，保留键级和原子属性。芳香分子的随书例子采用 JSON；不要仅凭某种交换格式中显示出了线条，就认定全部芳香信息已经保留。
\end{steps}

\ui{分子诊断}可以执行规则检查，当前图形界面结束后显示\ui{Molecule diagnostics completed.}，未展开详细问题表。图形操作中的检查按第~\ref{ch:modeling}~章\ref{sec:molecule-manual-check}节进行：先看组成和连接，再测量局部几何；完整MATLAB用户可从诊断返回数据读取逐项记录。

直接草绘时，按 S 进入工具，面板可指定元素、键级、形式电荷和杂化提示。从已有原子拖出可形成带键原子，拖到另一个已有原子可尝试建立键；空白处单击则放置原子。已有分子未选中连接对象时，工具不能替使用者决定哪一段应相连。出现碰撞提示时应移动草稿位置；按 Esc 退出工具，确认后续单击只用于选择。

\subsection{测量周期结构中的局部几何}\label{sec:periodic-measure-reference}

读取数值时，使用三维视图工具栏的\ui{Measurements}菜单，选择测量类型后按提示依次点选。顶部\ui{测量几何}当前完成后仅显示\ui{Geometry measurement completed.}，不在提示框中列出读数。查看器还提供晶胞参数、键统计、点到面距离、两平面夹角、配位数和最近邻。下表说明查看器中的选取顺序；这些工具回答不同的问题，不应只用两个原子的屏幕间隔代替结构检查。

\begin{table}[htbp]
\centering\small
\begin{tabularx}{\linewidth}{lX}
\toprule
任务 & 选取方式与读数 \\
\midrule
两点距离 & 按顺序取两个原子，距离用\AA。 \\
键角 & 取三个原子，第二个为顶点，角度用度。 \\
二面角 & 依次取四个原子，中间两个确定转轴，注意读数符号。 \\
点到面距离 & 先取待测点，再取三个定义平面的点；平面内三点不可共线。 \\
两平面夹角 & 依次取第一平面的三个点，再取第二平面的三个点。 \\
配位数或最近邻 & 选中心原子，查看当前邻接条件下的数目或距离列表。 \\
晶胞参数 & 从菜单直接读取当前晶胞，无须选原子。 \\
键统计 & 选两个原子以指定元素对，读取当前邻接算法找到的该元素对的键数、平均和最短最长距离。 \\
\bottomrule
\end{tabularx}
\caption{按问题选择几何测量方式。}
\end{table}

顶部\ui{测量几何}的返回数据在周期结构中使用邻近映像连接局部几何；查看器中的交互测量则按照实际点选的显示映像连线，选择哪一个边界映像会影响读数。最近邻列表会列出位点与映像偏移，读数时应一起查看。配位数还取决于邻接判定，不是对所有材料通用的固定化学常数。

要修改测得的几何量，可在相应测量编辑中输入目标值，也可使用顶部的\ui{设置距离}、\ui{设置键角}、\ui{设置二面角}。分子的\ui{atom}只动目标原子，\ui{subtree}与\ui{fragment}按连接寻找可移动部分；环不能被一根键分开时，移动范围会受限。在周期晶体中，这些精确几何命令当前只移动目标位点，不会自动把整段吸附分子当作分子子树。若要整体平移吸附物，应选择其全部位点，改用\ui{组分变换}。

\subsection{保存和连接常用片段}\label{sec:fragment-library-reference}
\index{Pianduanku@片段库}

片段草绘器把常见局部结构组织为可检索的库。片段的连接端口说明它怎样与宿主连接；有的片段适合单端连接，有的端口需要多个宿主位点，非共价端口则不提供这里的共价连接按钮。

\begin{steps}
  \item 在宿主分子中先选连接位点，再从\ui{分子建模}菜单打开\ui{片段草绘器}。搜索框可以缩小列表；选中条目后查看右侧三维预览、原子数、来源和标签。
  \item 在\ui{连接模式}中选择合适端口，阅读连接说明，核对宿主的选择数量。若\ui{连接片段}为灰色，先检查是否选了不适用于共价连接的端口。
  \item 单击\ui{连接片段}，检查新增部分的位置与宿主连接。内置条目的连接头已有定义，不宜仅凭预览中的球大小猜测连接原子。
  \item 应用后核对宿主与片段的连接、氢数及连接处的距离，再运行\ui{分子诊断}并保存项目。需要另一种连接方式时，撤销后重新选端口，避免在错误连接上继续叠加修改。
\end{steps}

要把自己的局部结构放入库，先让当前分子文档只包含准备复用的完整片段，检查连接和原子顺序，再打开\ui{保存用户片段}，填写名称。这个命令保存当前分子，不是仅保存屏幕中选亮的几个球。要替换已有同名条目时明确勾选覆盖，否则换一个名称。

片段草绘器中的\ui{导入片段库}与\ui{导出片段库}用于交换片段库 JSON。导出库文件便于在另一台电脑导入，与导出当前模型结构是两件事。导入后重新搜索条目并检查来源、端口和预览；完成连接后的模型仍应单独保存。

\section{处理周期晶体}

\subsection{改变晶胞而保留所需几何}\label{sec:lattice-transform-reference}

晶格矩阵的三行分别为 $\mathbf a,\mathbf b,\mathbf c$，三列为全局 $x,y,z$ 分量，单位为\AA。\ui{编辑晶格}中的矩阵是可编辑表格，应逐单元格填写；它不是把九个数一次输入一个文本框。

设编辑前后的晶格分别为 $A$、$A'$。\ui{保持笛卡尔坐标不变}决定了原子怎样适应新外框，其关系可直接从式\eqref{eq:row-lattice-coordinates}读出：
\begin{equation}
 \begin{aligned}
 \text{保留分数坐标：}\quad&\mathbf f'=\mathbf f,
       &\mathbf r'&=\mathbf f A',\\
 \text{保留笛卡尔位置：}\quad&\mathbf r'=\mathbf r,
       &\mathbf f'&=\mathbf r(A')^{-1}.
 \end{aligned}
 \label{eq:lattice-preservation-options}
\end{equation}
例如，立方晶胞边长由4改为8\,\AA，一个原子的原分数坐标为 $(0.5,0,0)$，原位置为 $(2,0,0)\,\text{\AA}$。不勾选保持笛卡尔位置时，它仍在晶胞的一半处，位置变为 $(4,0,0)\,\text{\AA}$；勾选时，位置保留为 $(2,0,0)\,\text{\AA}$，分数坐标改为 $(0.25,0,0)$。两种选择都可能有用，先决定要保留哪一种几何关系，再改表格。

\begin{table}[!htb]
\centering\small
\begin{tabularx}{\linewidth}{p{30mm}XX}
\toprule
操作 & 保留什么 & 重点检查什么\\
\midrule
编辑晶格，保留分数坐标 & $\mathbf f$和位点数 & 胞长、夹角及改变后的原子间距。\\
编辑晶格，保留笛卡尔位置 & $\mathbf r$和位点数 & 新分数坐标与周期边界接触。\\
构建超胞 & 原材料的周期排列 & $|\det S|$倍的体积和独立位点数。\\
旋转晶格 & 整体内部几何、位点数 & 晶格与原子是否一起转动。\\
交换坐标轴 & 笛卡尔几何、位点数 & 新基矢顺序与对应分数坐标。\\
镜像晶格 & 分数坐标、位点数 & 所选基矢取反及晶胞手性。\\
施加应变 & 分数坐标、位点数 & 三条基矢的长度及受影响的距离。\\
\bottomrule
\end{tabularx}
\caption{按希望保留的量选择晶格操作。表中的内部几何指原子间距离与角度；改变周期边界时，还需结合新晶胞检查邻近映像。}
\label{tab:round10-lattice-invariants}
\end{table}


\begin{steps}
  \item 保存原始晶胞，打开\ui{晶格编辑}中的\ui{编辑晶格}。先记下原胞长、夹角和代表性原子的笛卡尔位置。
  \item 逐格输入新矩阵。要让原子随晶格形变，保留分数坐标；要改变外框而保持当前笛卡尔位置，选择\ui{保持笛卡尔坐标不变}。例如，增大盒子装入孤立结构时，通常更关心后一种几何关系。
  \item 预览检查原子是否按预期移动，尤其检查跨边界的近邻。应用后核对晶胞体积与位置，再另存。
\end{steps}

\ui{重定义晶胞}直接使用新晶格矩阵并保留原子的笛卡尔位置；它不按新旧晶胞体积比复制原子。需要真正增加原子数时，应使用\ui{构建超胞}。

\subsubsection{反向一条晶格矢量}

\ui{镜像晶格}输入晶格矢量序号1、2或3，将相应矢量反向，其他两行不变。用一个对角矩阵 $D_i$ 表示这一操作：它的第 $i$ 个对角元素为 $-1$，另外两个为1，于是
\begin{equation}
 A'=D_iA,\qquad \mathbf f'=\mathbf f.
 \label{eq:mirror-lattice-row}
\end{equation}
例如选第1条轴，就是 $\mathbf a'=-\mathbf a$。分数坐标 $(u,v,w)$ 对应的位置由 $u\mathbf a+v\mathbf b+w\mathbf c$ 变成 $-u\mathbf a+v\mathbf b+w\mathbf c$。由于 $\det D_i=-1$，晶胞体积的大小不变，基矢的手性改变。

这与\ui{镜像原子}的作用不同。在一般斜晶胞中，反向 $\mathbf a$ 并不同时把 $\mathbf b$、$\mathbf c$ 关于某个笛卡尔平面反射，原子间的某些距离可能改变。若目的是将整个分子关于空间中的平面翻到另一侧，应使用前面的平面镜像公式；若要改晶格基矢方向，则按本式操作并核对原子位置。

\ui{旋转晶格}输入旋转角度和全局轴方向，会同时改变晶格与随晶格摆放的原子；其行矩阵关系为 $A'=AR^{\mathsf T}$，分数坐标保持不变。\ui{交换坐标轴}则同时重排晶格行与分数坐标分量：依次填2、1、3可交换 $\mathbf a$ 与 $\mathbf b$，同时交换 $u$ 与 $v$，因此笛卡尔位置不变。每次只做一种变化，按表\ref{tab:round10-lattice-invariants}核对后再保存。

\subsubsection{沿晶格方向施加应变}

\ui{施加应变}和\ui{结构应变}使用无量纲相对伸缩量。为突出作用方向，本书将三个输入记为 $\varepsilon_a,\varepsilon_b,\varepsilon_c$；它们依次缩放晶格矩阵的三行：
\begin{equation}
 A'=\begin{bmatrix}
 1+\varepsilon_a&0&0\\
 0&1+\varepsilon_b&0\\
 0&0&1+\varepsilon_c
 \end{bmatrix}A,\qquad \mathbf f'=\mathbf f.
 \label{eq:lattice-axial-strain}
\end{equation}
对话框中的三个分量仍以 $x,y,z$ 命名；在这个工具中应按第1、2、3条晶格矢量理解，斜胞时不能把它们当作沿全局 $x,y,z$ 分别拉伸。标量0.01使三行都乘1.01；输入\literalinput{0.01 0 0}只使 $\mathbf a$ 乘1.01。0.01表示伸长1\%，不是增加 $0.01\,\text{\AA}$，也不是将胞长直接设为1.01\,\AA。

仍用第\ref{ch:modeling}章的手算斜胞，$\mathbf a=(4,0,0)$、$\mathbf b=(1,3,0)$、$\mathbf c=(0,0,6)$，单位为\AA。施加\literalinput{0.01 0 0}后，只有 $\mathbf a$ 变为 $(4.04,0,0)\,\text{\AA}$；$\mathbf b$ 的 $x$ 分量仍为1\,\AA。分数坐标 $(0.5,0.5,0)$ 不变，笛卡尔位置则由 $(2.5,1.5,0)$ 变为 $(2.52,1.5,0)\,\text{\AA}$。这个结果也说明，它不是把所有原子的全局 $x$ 坐标统一乘1.01。

先保存未变形模型，再填写应变、预览并应用。按上述小比例检查基矢长度和代表性原子的位置，最后另存结果。当前这两个工具只提供三条晶格矢量的伸缩，不提供剪切应变；施加形变后若需要平衡结构，应另行设置相应的优化计算。


\subsection{建立正交超胞}\label{sec:orthogonal-cell-reference}

有些后续操作希望晶胞角接近 $90^\circ$。\ui{晶胞正交化}搜索与原晶格相容的整数超胞，不是任意把斜晶胞拉成直角。模型可能变大，原子数和体积应一起检查。

\begin{steps}
  \item 从原始小晶胞开始，在\ui{超胞与晶格}中打开\ui{晶胞正交化}。
  \item 填写最小、最大胞长，单位为\AA；最大值必须大于最小值。再给定原子数上限和角度容差，容差以度表示。界面会据当前晶格提供起始长度范围，默认原子数上限5000、直角容差 $0.1^\circ$。
  \item 先用适合自己计算规模的原子上限预览。完成后检查三个角与 $90^\circ$ 的差、实际胞长以及位点数，而不是只看外框在当前相机下像不像矩形。
  \item 若无候选，依次扩大最大胞长、适当增加原子上限或放宽角度容差，保留每次使用的条件。取消失败或不合适的预览，回到原模型。
  \item 应用后保存为有区别的文件名，并记下正交化后的独立位点数。计算输入应引用这个新模型。
\end{steps}

\subsection{检查对称性并选择标准晶胞}\label{sec:symmetry-reference}

\ui{分析对称性}根据位置和晶格判断空间群。它不修改结构；\ui{原胞}与\ui{惯用胞}则会生成标准表达，位点数和坐标可能改变。名称相近的\ui{转为 P1}用于去除对称性约束，适合开始局部编辑，不能用它代替原胞约化。

\begin{steps}
  \item 在\ui{对称性工具}中选择\ui{分析对称性}，从位置容差 $0.01\,\text{\AA}$、角度容差 $5^\circ$ 开始查看。读取空间群符号、Hall符号和国际编号，保存到自己的建模记录。
  \item 对有轻微数值偏差的输入，可逐次比较更细或更粗的位置容差。容差是在判断近似等价，不是一次结构弛豫；不能仅为得到想要的空间群而无限放宽。
  \item 要比较原胞与惯用胞，从同一原始模型分别执行两项构建，各自另存。比较化学式、体积、位点数及元素比例，确认它们表达相同的周期材料。
  \item 要对晶格本身作空间观察，打开\ui{Wigner-Seitz 晶胞}。结果显示由晶格定义的多面体；它不把原子裁成一个团簇，也不改变当前模型位点。需要结构文件时仍保存原晶体。
\end{steps}

开始局部缺陷编辑前，可再练习一次 P1 转换。打开按空间群构建后保留的\file{Cu-fcc}模型，在\ui{对称性工具}中选择\ui{转为 P1}；这项命令没有额外数值参数。执行后核对四个 Cu 位点和坐标仍与原模型一致，再另存为\file{Cu-P1}。如果随后只改一个局部位置，应在这个副本上继续，保留原有空间群模型供比较。P1 表示不再施加额外的空间群约束，不会自动让原来规则的排列变得不规则。

\subsection{明确哪些原子允许移动}\label{sec:fixed-sites-reference}

\ui{固定原子}为模型写入选择性移动标记。它用于说明后续结构弛豫的约束，不能仅靠三维视图中某个原子看起来不动来判断约束是否存在。

\begin{steps}
  \item 选中准备固定的全部位点，打开\ui{原子编辑}中的\ui{固定原子}，核对原子序号。
  \item 若薄片带有几何层信息，界面还会出现固定底层数。要固定最低若干层时输入层数；要严格使用自己的位点列表，将底层数设为0。
  \item 应用后保存项目，并导出原生结构JSON。重新打开该结构，核对位点数与准备固定的序号。若需要直接检查记录，可用文本编辑器只读查看这些位点的\file{selective_dynamics}：三个false表示三个方向都固定，三个true表示都可移动。
\end{steps}

该命令本次选中的位点设为三个方向均固定，其余位点设为可移动。因此它不是每按一次就在原有固定集合上追加若干原子的开关；若希望固定两组，应先把两组合成完整选择再执行。保存选择集可以减少重复选取时的遗漏。

在分子建模的\ui{几何优化}窗口，约束由\ui{固定原子索引}填写；例如输入\literalinput{1 2}可固定两个指定原子，再预览优化结果。本节的晶体固定标记用于记录位点约束，不能用它替代计算任务本身的约束设置。第~\ref{ch:workflows}~章的氢分子优化不含固定原子；需要带约束计算时，先按所用计算方法确认约束输入，再检查优化前后固定原子的坐标。

\subsection{建立空位以外的点缺陷}\label{sec:other-defects-reference}

\ui{缺陷与合金}\,$\rightarrow$\,\ui{创建点缺陷}还包含替位、间隙和反位。它们改变原子的方式不同，应从预先保存的宿主副本分别建立，避免把几个缺陷无意叠在一起。

\begin{steps}
  \item 建立替位缺陷时，选择\ui{替位}，指定一个原子序号与新元素。应用后位点数应保持不变，组成中相应元素的数目发生一增一减。
  \item 建立间隙缺陷时，选择\ui{间隙}，输入元素，明确选择分数或笛卡尔坐标，再填位置。预览检查最近邻接触，应用后位点数应增加一个。不能因为几何中心容易输入，就认为晶胞中心一定是合适间隙。
  \item 建立反位缺陷时，选择\ui{反位}，填写两个位点序号。这会交换二者的物种，保留位置与总组成；两个相同元素互换不会产生有意义的化学反位。
  \item 每个结果单独命名，记录宿主、超胞大小与缺陷位置。改变原胞表达或排序后，重新核对序号。
\end{steps}

\ui{枚举点缺陷}用于空位或替位候选，输入允许考虑的位点范围、替换元素、参考电荷和对称容差。结果表中的\ui{Degeneracy}表示等价候选数，\ui{NetCharge}与输入的参考电荷有关，最短接触用于几何检查；它们不等于缺陷形成能或已经求得的电子电荷态。应用枚举后，各缺陷结构自动加入项目浏览器的结构分组，并打开相应文档；结果表用于查看，不在表内切换模型。到项目浏览器逐个打开候选，核对组成与位点数，再保存项目。后续比较采用一致的宿主与数值条件。

\subsection{生成显式占位的合金模型}\label{sec:sqs-reference}

SQS（特殊准随机结构）在有限晶胞中安排具体的元素位置，使邻近元素配对等统计量尽量接近随机混合。它把输入的混合占位转成每个位点由一种元素完整占据的模型，便于继续计算；这不表示材料已经变为化学有序相。命令读取的是输入结构中的元素与占位比例，不会根据纯元素晶体自动决定合金成分。

下面的四位点铜镍实例给出输入文件、倍率、搜索时限和结果核对。换用其他组成时，先选择能容纳整数原子个数的晶胞；例如等比例两组分需要能分成相等的整数数量。保存时同时保留混合占位输入和生成后的具体排列。

\subsection{示例：建立四位点铜镍SQS}\label{sec:sqs-cu-ni-example}
\index{SQS@SQS 模型}

随书\file{sqs-cu-ni-input.json}提供一个可以直接开始的输入：边长为 $3.61\,\text{\AA}$ 的面心立方晶胞，含4个独立位点。每个位点均记录Cu占位0.5、Ni占位0.5，总组成相当于Cu$_2$Ni$_2$。这里的混合占位描述统计上的组分分布，并不是在同一位置放置两个实际原子。SQS将这些比例转成明确的原子种类，使每个位点由一种元素完整占据。

本例的晶格常数仅作教学初值。四位点模型很小，适合看清输入与输出的区别；它不用于判断铜镍合金的平衡晶格常数，也不凭这一小晶胞代表宏观无序材料的全部关联。

\begin{table}[htbp]
\centering\small
\begin{tabularx}{\linewidth}{lX}
\toprule
项目 & 本例设置与检查 \\
\midrule
输入文件 & \file{sqs-cu-ni-input.json}；另附同名CIF。\\
晶格 & $a=b=c=3.61\,\text{\AA}$，三个夹角均为 $90^\circ$。\\
独立位点 & 4；分数坐标为 $(0,0,0)$、$(0,\frac12,\frac12)$、$(\frac12,0,\frac12)$、$(\frac12,\frac12,0)$。\\
位点占位 & 每个位点Cu 0.5、Ni 0.5；两种占位之和为1。\\
SQS 超胞倍率 & 在文本框中输入\literalinput{1}，保持本例四位点晶胞。\\
SQS 搜索时限 & 输入1秒；本例进入小规模精确枚举。\\
预期结果 & Cu 2个、Ni 2个，共4个位点，每个位点占位为1。\\
\bottomrule
\end{tabularx}
\caption{四位点铜镍SQS练习。}
\end{table}

\begin{steps}
  \item 将\file{sqs-cu-ni-input.json}导入项目，打开其三维文档。Online中先上传文件，再从云端位置导入。将文档命名为\file{CuNi-mixed}，保存起始项目，以便随时回到混合占位输入。
  \item 查看独立位点数、晶胞和原子信息。位点数应为4，位点组分应同时包含Cu和Ni，各为0.5。若查看输入JSON，每个\file{sites}条目中的\file{species}都有两个元素及各自的\file{occu}值；这比仅凭显示颜色判断组成更可靠。
  \item 保持该晶体文档激活，从\ui{晶体建模}\,$\rightarrow$\,\ui{缺陷与合金}打开\ui{生成 SQS 模型}。在\ui{SQS 超胞倍率}中输入\literalinput{1}，在\ui{SQS 搜索时限（秒）}中输入1。
  \item 单击\ui{预览}，查看四个位点是否已分配成两种明确的元素，确认晶胞大小和位点数没有增加，再\ui{应用}。本例保留原有晶格和位置，改变的是各个位点的物种占据。
  \item 检查结果组成：应为Cu$_2$Ni$_2$，位点数仍为4，每个位点只含一种元素、占位为1。测量最近邻距离，应约为 $3.61/\sqrt{2}=2.552655\,\text{\AA}$。元素分配可以通过晶胞对称操作得到等价表达，检查时以组成、占位和周期几何为准。
  \item 按前述结构导出方法导出原生JSON，命名为\file{sqs-cu-ni-ordered.json}，保留搜索记录。文件名中的\file{ordered}表示每个位点已有明确元素，不表示生成了化学长程有序相。
  \item 将结果文档命名为\file{CuNi-SQS-4}并保存项目；再导出CIF，便于交换晶胞和明确的元素占位。Online中下载导出文件，并保留输入JSON、结果JSON和项目副本。CIF用于结构交换，SQS搜索记录以原生JSON为准。
\end{steps}

\subsubsection*{查看详细搜索记录}

用文本编辑器打开导出的JSON，只查看\file{properties.sqs}中的记录，无需修改文件。\file{search_mode}表示搜索方式，\file{configuration_count}表示候选排列数，\file{objective_function}是关联匹配的目标值。

随书脚本本次得到的\file{search_mode}为\literalinput{exact-enumeration}，\file{configuration_count}为6，\file{objective_function}约为0.1666667。六种排列来自四个位点中选择两个放Cu，其余放Ni。本例六对不同位点的周期最短距离相同；Cu--Cu与Ni--Ni各占其中一对，有限晶胞不能完全再现独立随机占位的配对比例，因此目标值并不为零。这个无量纲目标值衡量当前关联函数与随机占位目标的差别，不是吸附能、总能量或缺陷形成能。

搜索时限用于控制搜索预算，不是全部建模操作的倒计时。本例完成精确枚举、等价结构比较及结果整理约用4.4秒，即使填写1秒，也应等待操作结束；不同计算机上的时间会有所不同。后续扩大模型时，还应同时考虑位点数与候选规模。

\begin{notebox}{倍率与混合占位}
本命令的单个整数倍率表示晶胞体积倍数，不是三个方向各乘同一个数。要明确三个方向的重复数，在同一倍率文本框中输入例如\literalinput{2 2 2}；它会把本例扩为32个位点。比较模型前应记录实际晶胞、组成和搜索条件。混合占位输入用于建模；电子结构计算应使用生成并检查过的显式占位结果。导入混合占位时，计算对象可能给出分数占位不能直接表示的提示，不能把转换出的单一元素显示当成SQS已经完成。
\end{notebox}

随书\file{sqs-cu-ni-ordered.json}和\file{sqs-cu-ni-ordered.cif}提供本次参考结果；\file{sqs_example.m}可在MATLAB无JVM模式下重新生成输入、执行同一建模命令并回读核对文件。这组参考文件用于核对结构与占位；电子结构计算可从生成的模型继续安排。

% SCREENSHOT: CuNi混合占位输入的4sites与SQS倍率1/时限1秒参数；应用后的Cu2Ni2与导出JSON搜索记录，须由实机采集。


\subsection{选择低维结构的构建方式}\label{sec:nano-reference}

第~\ref{sec:nanostructure-examples}~节已用石墨烯和铜完成五类构建。这里按希望保留的边界选择入口，各例参数见该节。

\begin{table}[htbp]
\centering
\begin{tabularx}{\linewidth}{lX}
\toprule
需要的模型 & 命令与关键区别 \\
\midrule
有限宽度的片层 & \ui{截取纳米带}；按宽度、长度重复数裁切，沿带长周期。\\
卷曲的片层 & \ui{生成纳米管}；按手性指数与轴向重复数构建，沿管轴周期。\\
有限横截面的晶体 & \ui{截取纳米线}；按晶格方向与半径截取，保留轴向周期。\\
三个方向均有限的团簇 & \ui{量子点 / 纳米空洞}的\ui{量子点}模式；保留球内原子。\\
周期块体中的球形孔洞 & 同一窗口的\ui{纳米空洞}模式；去掉球内原子，保留外围周期宿主。\\
\bottomrule
\end{tabularx}
\caption{按边界选择纳米结构命令。}
\end{table}

线、带和管的重复数不是以\AA 为单位的总尺寸。比较不同模型时，以实际晶胞、位点数和周期方向为准；需要钝化时，先保存未钝化输入，再检查新增原子的组成与边界配位。

\section{建立软物质模型}

\subsection{构建带有支链的高分子}\label{sec:branched-polymer-reference}

支化构建从一条主链出发，按重复单元间隔连接片段，适合学习主链与支链的拓扑关系。重复单元名称来自重复单元库，支链名称来自片段库，两者应分别核对。

\begin{steps}
  \item 先导入随书\file{user-repeat-ethane-input.json}并激活分子文档，从\ui{分子建模 → 软物质 → 聚合物建模}打开\ui{构建支化聚合物}。保留输入副本；构建结果会替换活动文档中的分子。先用短链练习，例如主链重复单元\literalinput{PE}、重复数8、支链\literalinput{Methyl}、支化间隔3、种子1。
  \item 查看预计原子规模，给出适合练习的原子上限。支化间隔按重复单元计数，不是按笛卡尔长度或主链上每一个原子计数。
  \item 预览并检查支链是否连在预期主链单元附近，再确认主链的两端与连接。若需改变支链密度，回到原始短链调整间隔，不要在已经支化的结果上反复执行构建。
  \item 应用后检查组成、支链连接和局部距离，再运行\ui{分子诊断}并保存原生模型。密度与装填批次属于装填任务，不能把本次主链构建的完成当成已经得到指定密度的材料。
\end{steps}

\ui{构建均聚物}可用于不含支链的对照：选一个重复单元、重复数、独立链数、端基与初始构象，固定种子后构建。独立链数大于1表示生成多条链，不是自动建立分子间共价键。\ui{伸展}、\ui{折线}与\ui{无规线团}是初始几何选择；保存后仍应检查链内接触及后续计算是否需要进一步整理。

\subsection{建立基础树枝状骨架}\label{sec:dendrimer-reference}

\ui{构建基础树枝状聚合物}按代数与分支数展开骨架，适合观察拓扑层次。核心官能度决定第一代分支数，分支因子决定后续每个分支继续产生多少分支。界面的骨架元素和键长用于初始几何，不能由“树枝状”这个名称推定它已经代表某种具体合成材料。

\begin{steps}
  \item 导入随书\file{user-repeat-ethane-input.json}并激活，保留输入副本。从\ui{分子建模 → 软物质 → 聚合物建模}打开\ui{构建基础树枝状聚合物}，先用代数2、核心官能度3、分支因子2、元素C、骨架键长 $1.54\,\text{\AA}$、种子1作小模型练习。
  \item 查看预计规模，设置原子数上限。增加代数会使分支快速增多，先看小模型的连接再扩展。
  \item 预览时找到中心，顺着连接查看第一代和第二代。检查分支是否意外穿过其他分支，查看最短接触。
  \item 应用后检查组成、分支连接和局部接触，运行\ui{分子诊断}后保存项目和原生结构。几何整理或后续弛豫应另行进行；构建成功仅说明已经得到这一层次的初始骨架。
\end{steps}

\subsection{保存自定义重复单元}\label{sec:repeat-unit-reference}

重复单元除了原子和键，还需要两个不同的连接端。头原子用于连接前一个单元，尾原子用于连接后一个单元；离去原子在连接前移除。下面用乙烷中的两个碳定义一个简单骨架，再连接成二聚体，查看这些字段怎样影响实际结构。

随书\file{user-repeat-ethane-input.json}含两个C和六个H，共8个原子、7条单键。C为1、2号，H为3至8号。输入中的C--C长度为$1.52\,\text{\AA}$。本例移除六个H后保存碳骨架，构建链时再由工具统一补氢；这是建立单元的几何操作，不用于表示乙烷的聚合反应。

\begin{table}[htbp]
\centering
\begin{tabularx}{\linewidth}{lX}
\toprule
操作字段 & 本例填写 \\
\midrule
用户重复单元名称 & \literalinput{MANUAL_CC}。\\
头原子、尾原子 & 分别为1、2。\\
离去原子序号 & 在文本框输入\literalinput{3 4 5 6 7 8}。\\
均聚物重复单元 & 输入保存的名称\literalinput{MANUAL_CC}。\\
重复次数、链数 & 分别设为2、1。\\
初始构象、种子 & 选择\ui{伸展}，种子7。\\
首尾端基 & 均保持\literalinput{none}；其余未满足价键由氢补足。\\
原子数上限 & 1000；足以容纳本次小模型。\\
\bottomrule
\end{tabularx}
\caption{从自定义碳骨架建立二聚体。}
\end{table}

\begin{steps}
  \item 导入\file{user-repeat-ethane-input.json}，打开分子文档。逐个查看前两个原子，确认1、2号是碳；再核对3至8号均为氢。保留此输入文件，不在建立重复单元前重新排序。
  \item 从\ui{分子建模 → 软物质 → 聚合物建模}打开\ui{保存用户重复单元}，按表输入名称、头尾及离去原子序号。头尾必须是不同的现有原子，且不能同时列为离去原子。
  \item 首次保存不勾选覆盖。若已有同名教学条目，确认是本例旧记录后再覆盖，或改用另一个名称并在后续保持一致。保存后当前乙烷仍是8个原子；保存单元本身不修改来源分子。
  \item 仍在这个分子文档中打开\ui{构建均聚物}。填入保存的名称，按表设置重复次数、链数、构象、端基和种子，头尾连接方式保持默认。预览后应用，检查生成的完整链。
  \item 结果应有4个C、10个H，共14个原子、13条单键，两个单元之间新增一条C--C连接。旋转检查主链与末端氢，对可疑接触测量距离，再运行\ui{分子诊断}。随书\file{user-repeat-CC-dimer.json}的价态与接触检查数据由配套脚本保留，可用于对照。
  \item 将结果命名为\file{MANUAL-CC-dimer}，保存项目并导出原生结构JSON。保存输入乙烷、重复单元参数和二聚体三个阶段；只交换二聚体文件不会把自定义单元自动注册到另一台电脑的库中。
\end{steps}

要查看详细来源，可用文本编辑器只读打开输入或结果JSON，比较元素、坐标和连接；完成本例不需要手工修改文件。随书\file{user_repeat_unit_example.m}重复执行相同的保存与构建，并将教学库单独写到\file{user-repeat-unit-library.json}，不覆盖个人重复单元库。参考二聚体的组成、连接和诊断结果，以及导出后的读回结果，均可用配套核对脚本检查。继续建立更长链时，从相同单元开始增加重复次数，再检查完整链内的接触。

\subsection{将一种分子装入受限区域}\label{sec:confined-packing-reference}

\ui{构建无定形晶胞}以当前分子为组分，按分子数和密度放入周期盒；\ui{构建受限分子层}进一步限定可用区域。密度单位为 $\mathrm{g/cm^3}$，最小接触用\AA。装填得到的是初始位置，后续物理平衡需要单独完成。

\begin{steps}
  \item 打开连接与几何已检查的分子。先保存组分原件，再从\ui{无定形建模}选择单组分装填或受限分子层。
  \item 输入分子数、目标密度、最小接触、种子与原子上限。首次练习使用较少分子和较低初始密度，先确认组分及边界的处理方式。
  \item 对受限分子层，选择限制轴1、2或3，再填写区间上下限。例如 \literalinput{0.2 0.8} 表示该轴盒长的20\%至80\%；当两个数都在0与1之间时按比例解释，不能用它表达 $0.2$ 到 $0.8\,\text{\AA}$。
  \item 预览时从限制轴方向的侧面检查分子是否落在区域内，确认没有把层的空白部分误当作可装填体积。
  \item 应用后核对分子数、总原子数、盒长以及跨周期接触，另存结果与输入条件。
\end{steps}

\ui{密度递增阶段（可选）}可填写一串正的、非递减的密度，最后一个值必须等于目标密度；留空则只用目标密度。当前密度渐增不能和受限区域等限制任意叠加，做受限层练习时先留空。\ui{每批装填分子数}控制装填批次，不改变总分子数。需要按两种分子组成装填时，使用前面的\ui{装填混合物}教程；摩尔分数和质量分数模式还需给出总分子数，最终应检查换算后的整数数量。

\subsection{向已有盒子装填分子}\label{sec:host-packing-reference}

已有宿主时，\ui{装入已有晶胞}保留当前晶胞与宿主原子，把另一个项目分子加入盒中。它不根据目标密度重新建立一个立方盒。\ui{绕纳米粒子装填}在此基础上，为当前宿主周围留出指定间隙。

\begin{steps}
  \item 把宿主和待装填分子导入同一项目，分别命名。激活具有正交晶胞的宿主文档，在\ui{宿主装填}中选择需要的命令；斜晶胞需先整理为适合装填的盒子。
  \item 在\ui{待装填分子}中选择组分，填写分子数、最小接触、种子、每批分子数和原子上限。这里没有通过密度重新计算盒长的步骤，应先保证宿主盒中有足够空间。
  \item 绕纳米粒子装填时再填写粒子间隙，单位为\AA。排除区域以宿主平均位置为中心、覆盖宿主最外原子的球形范围再加间隙；它不逐个追随任意凹凸表面。
  \item 预览并从不同方向检查宿主保留、分子位置与跨周期边界的接触。空间不足时，回到原宿主扩大适当的空白盒子，或减少装填数量后重新开始。
  \item 完成后核对新增分子数与组成，将宿主、组分原件和装填结果分别保留；保存项目及原生结构。
\end{steps}

\section{完善表面模型}

\subsection{钝化表面}\label{sec:surface-finishing-reference}

第~\ref{ch:advanced}~章的吸附练习介绍了候选位点与O--H放置。这里处理薄片上仍需封端的边界；从已保存的薄片副本开始，使钝化和后续溶剂处理分别留下记录。

\begin{steps}
  \item 若表面有未饱和配位，需要添加封端原子，打开\ui{表面钝化}，选择元素、目标键长、处理上表面/下表面/两面，以及表面探测深度。
  \item 目标配位数为0时由工具推断；已知配位目标时明确输入。最小接触用来排除过近位置。预览中检查两侧添加是否符合目的，并查看原有原子有无被错误识别为表面。
  \item 应用后比较新增元素数与边界配位，单独保存钝化结果，再决定是否继续装填溶剂。
\end{steps}

\subsection{在薄片上方加入溶剂}

\ui{添加溶剂层}需要项目中已有一个完整溶剂分子。保持薄片文档激活，在命令对话框的\ui{溶剂分子/结构}下拉菜单中选择它，再设置分子数、层厚、表面间隔、剩余真空、最小接触、种子和超时秒数。层厚、间隔、真空和接触都用\AA；分子数为整数。先用少量分子查看它们位于薄片哪一侧，核对溶剂内部分子几何、层内接触和剩余真空后再扩大。该命令的完成不表示溶剂层已经达到指定温度或平衡密度。

\section{复用建模过程}

\subsection{保存常用的建模参数}\label{sec:modeling-preset-reference}
\index{Canshuyushe@参数预设}

参数预设属于某一个命令。它保存当前填写的条件，不保存一个新的模型，也不把同样参数自动换算为适合另一材料的数值。

\begin{steps}
  \item 打开常用建模命令，填写并核对参数。可以先预览检查效果；要保存一组待使用的条件，也可在应用前保存。
  \item 单击对话框顶部的\ui{保存预设}，输入说明用途的名称，例如“表面练习用厚度”。保存后确认该名称出现在\ui{预设}列表中。
  \item 关闭并重新打开同一命令，选中名称，再单击\ui{加载}。逐项查看参数是否符合当前模型，特别是位点序号、项目结构引用和长度。
  \item 必要时改动参数，再预览、应用并保存结果模型。载入预设本身不代表建模已经提交。
\end{steps}

对话框底部\ui{重置}恢复这个对话框的初始填写值，同时清除临时预览；顶部历史区的\ui{重置}则恢复当前结构草稿的基准，两者作用不同。对话框中的\ui{帮助信息}可展开查看参数解释，再收起继续填写。

\subsection{录制可追溯的建模过程}\label{sec:modeling-recipe-reference}
\index{Jianmopeifang@建模配方}

Recipe记录一系列已经提交的建模命令、参数、随机种子和前后结构的校验信息。它适合保留“这个结果怎样得到”的过程；原始模型与最终结构仍应同时保存。

\begin{steps}
  \item 打开并保存一份明确命名的起始模型。激活其三维文档，在\ui{历史}中单击\ui{录制}。
  \item 执行需要记录的建模操作。先预览、检查，再应用；取消的试验不应当作已经完成的结构步骤。
  \item 这次过程结束后，单击\ui{保存 Recipe}。录制随之停止，软件按时间生成名称，并在状态信息中给出保存位置。
  \item 打开\ui{历史}\,$\rightarrow$\,\ui{预设与模板}，单击\ui{刷新}，在类型为\ui{操作配方}（Recipe）的条目中查找本次记录。保存模型结果和项目，并把起始输入一同归档。
\end{steps}

录制针对当前结构文档中受支持的建模提交，不能当作录屏或全部项目活动的日志。视图变化、选择集管理和单纯查看分析不构成同样的结构命令记录；从分析结果窗口应用候选后，也应把最终模型另存。当前\ui{预设与模板}窗口提供查看已有条目和新建模板；没有在选中一条Recipe后直接重放到任意当前结构的按钮。需要重放时，MATLAB用户可采用建模API，并使用与记录匹配的起始结构；普通图形操作只需把记录保留为可复查的过程文件。

\subsection{建立建模项目模板}\label{sec:modeling-template-reference}

\ui{预设与模板}把命令预设、Recipe和建模项目模板放在同一列表中，列出类型、名称与详情。这里的建模项目模板记录所用Recipe名称及输入输出格式，与工作流页面保存的 \file{.wf} 计算模板用途不同。

\begin{steps}
  \item 在\ui{历史}中打开\ui{预设与模板}，先单击\ui{刷新}，确认准备引用的Recipe存在。
  \item 单击\ui{新建项目模板}，填写名称、Recipe名称、输入格式和输出格式。默认\literalinput{auto}表示由输入识别格式，\literalinput{same}表示输出沿用输入格式；它们是格式选择，不是文件路径。
  \item 保存后在列表核对类型、名称和详情中的输入到输出关系。若条目标为无效，按详情检查相应记录，不要把它当成已经生成好的项目。
  \item 将模板和它引用的Recipe一起保留。实际模型文件另行保存；创建模板不会自动打开一份新工程或执行建模。
\end{steps}

\subsection{查看已有的批处理任务}\label{sec:modeling-jobs-reference}

\ui{建模任务}查看已经存储的建模文件批处理任务，与电子结构工作流的SCF运行列表不同。窗口显示编号、状态、已完成数、请求数、成功数、是否可恢复以及错误信息；一个批次完成，也可能存在个别文件失败。

\begin{steps}
  \item 打开\ui{历史}\,$\rightarrow$\,\ui{建模任务}，单击\ui{刷新}，按任务编号找到需要查看的批次。
  \item 比较请求数、已完成数与成功数，再读错误栏。不要只依据状态栏的一句“完成”判断每个输入都得到了输出。
  \item 选中可恢复的条目，再单击\ui{恢复任务}，等待状态更新并刷新列表。若按钮为灰色，先确认该任务当前是否标为可恢复。
  \item 不准备继续的可恢复任务可单击\ui{取消任务}请求停止，随后刷新确认状态。取消不是删除已经写出的结果，输出文件仍需按任务记录检查。
  \item 到任务原有输出目录检查结果结构与报告，保留失败条目的输入，完成正常的数据归档。
\end{steps}

这个窗口用于查看与恢复已有任务。需要创建文件批处理时，完整MATLAB用户可使用建模API；图形界面用户按上述方法检查已有任务。

\subsection{恢复尚未保存的结构编辑}\label{sec:modeling-draft-reference}

撤销与重做针对当前结构草稿；历史区的\ui{重置}恢复草稿基准，保存后基准随模型更新。要长期比较两个模型，应分别保存命名结果。

重新打开文档时若提示恢复草稿，核对对应结构后再恢复，并检查原子数、组成、晶胞和最后编辑位置。另存结果；恢复不会重建中断前的全部撤销步骤。若不恢复，先确认原保存文件就是需要的版本，再清理草稿。

在三维视图中，空格键重新居中，I切换简洁视图；Ctrl/Cmd配合加号、减号或0调整界面缩放。这些操作不改变原子坐标。输入参数后，将焦点移回视图再使用快捷键；按\ui{?}查看快捷键，或从\ui{历史 → 建模指南}查阅帮助。
